%% Bibliography and Glossaries Configuration

% csquotes: Context-sensitive quotation marks (required with babel + biblatex)
\usepackage{csquotes}

% biblatex with IEEE style: Modern bibliography management
\usepackage[
    style=ieee,
    backend=biber,
    natbib=true,
    hyperref=true
]{biblatex}

% Load bibliography database
\addbibresource{config/quellen.bib}

% glossaries: Acronyms and symbols lists
% Options: toc (in TOC), section=chapter (separate chapters), nonumberlist (no numbers)
% acronym (extra acronym list), symbols (extra symbol list)
\usepackage[toc, section=chapter, nonumberlist, acronym, symbols]{glossaries}

% Style for glossaries: long-short format with hyperlinks
\setglossarystyle{long}

% Initialize glossaries (required for all entries)
\makeglossaries

% ===== ACRONYMS / ABKÜRZUNGEN =====
% Format: \newacronym{key}{short}{long}
% Usage in text: \gls{key} = long (short) on first use, then (short)
% Usage after first: \gls{key} = (short)
\newacronym{ide}{IDE}{Integrated Development Environment}
\newacronym{ocxo}{OCXO}{Oven Controlled Crystal Oscillator}
\newacronym{ic}{IC}{Integrated Circuits (Integrierte Schaltkreise)}
\newacronym{ntc}{NTC}{Negative Temperature Coefficient (Heißleiter)}
\newacronym{gpio}{GPIO}{General Purpose Input/Output}
\newacronym{usbc}{USB-C}{Universal Serial Bus Type-C}
\newacronym{ufp}{UFP}{Upstream Facing Port}
\newacronym{adc}{ADC}{Analog-to-Digital Converter (Analog-Digital-Wandler)}
\newacronym{rms}{RMS}{Root Mean Square (Quadratischer Mittelwert)}
\newacronym{cmos}{CMOS}{Complementary Metal-Oxide-Semiconductor}

% ===== SYMBOLS / SYMBOLVERZEICHNIS =====
% Format: \newglossaryentry{key}{type=symbols, name=Symbol, description=..., symbol=Unit}
% Usage in text: \gls{voltage} or \glssymbol{voltage} for unit only
\newglossaryentry{voltage}{
    type=symbols,
    name=$U$,
    description={Elektrische Spannung; Potentialdifferenz zwischen zwei Punkten},
    symbol={V}
}

\newglossaryentry{resistance}{
    type=symbols,
    name=$R$,
    description={Elektrischer Widerstand; Widerstand gegen Stromfluss},
    symbol={$\Omega$}
}

\newglossaryentry{current}{
    type=symbols,
    name=$I$,
    description={Elektrische Stromstärke; Menge der Ladung pro Zeiteinheit},
    symbol={A}
}

\newglossaryentry{power}{
    type=symbols,
    name=$P$,
    description={Elektrische Leistung; Energieumsatz pro Zeiteinheit},
    symbol={W}
}

\newglossaryentry{frequency}{
    type=symbols,
    name=$f$,
    description={Frequenz; Anzahl der Schwingungen pro Sekunde},
    symbol={Hz}
}
